\documentclass[12pt, twocolumn, landscape, a4paper]{article}

\usepackage[margin=1cm]{geometry}
\usepackage[utf8]{inputenc}
\usepackage[spanish,es-noshorthands]{babel}
\usepackage{amssymb, amsmath, amsbsy}
\usepackage{tasks}
\usepackage{enumerate}
\usepackage{mathpazo}
\usepackage[pdftex]{hyperref}
\hypersetup{pdftitle={Ejemplos con conjuntos},pdfauthor={Luis Carrillo Gutiérrez}}

\renewcommand{\familydefault}{\sfdefault}
\pagestyle{empty}

\begin{document}
\paragraph*{Conjuntos}

\begin{itemize}
\item{Sea $A$ un conjunto con dos elementos y sea $B$ un conjunto con tres elementos, el número de elementos de $n\big[P(A)\big]\times n[P(B)]$ es :
\begin{tasks}(5)
\task{12}\task{24}\task{32}\task{48}\task{64}
\end{tasks}
}
\item{Si $A$ y $B$ son dos conjuntos incluídos en el conjunto universo, tales que $n(A) = 12$; \quad$n(B) = 16$\quad y \quad$n(A \cap B^{\,c}) = 7$. \quad Calcular $n(A\triangle B)$

\begin{tasks}(5)
\task{10}\task{11}\task{17}\task{18}\task{19}
\end{tasks}} 
\item{Determinar el valor de verdad o falsedad de las siguientes proposiciones :
\begin{enumerate}[I.]
\item{$\{2\} \cup \{\varnothing\} = \{2\}$} %F
\item{Si $\{\varnothing\} \in B$, entonces $\{\{\varnothing\}\} \subset B$, donde $P(B)$ es potencia de B }
\item{Si $A = \{\varnothing;\,\{\varnothing\};\,\{\varnothing;\,\{\varnothing\}\}\}$ entonces el conjunto $P(A)$ tiene 8 elementos}
\end{enumerate}  
\begin{tasks}(5)
\task{VVV}\task{VFV}\task{FVV}\task{FFV}\task{VFF}
\end{tasks}} 
\item{Un club deportivo tiene $68$ jugadores, de los cuáles $48$ practican el fútbol, $25$ el basquet y $30$ el béisbol. Si solo $6$ jugadores practican los tres deportes. ¿Cuántos jugadores practican exactamente un deporte? 
\begin{tasks}(5)
\task{30}\task{36}\task{39}\task{41}\task{43}
\end{tasks}}
\item{Sean los conjuntos : $A = \{2;\,3;\,8\}$\quad y \quad$B = \{1;\,2;\,7\}$ con los siguientes enunciados :
\begin{enumerate}[I.]
\item{$\exists\,x \in A, \forall\,y \in B\;/\; x + y \geq 9$}
\item{$\exists\,x \in A, \exists\,y \in B\;/\; x + y = 4$}
\item{$\forall\,x \in A, \forall\,y \in B\;/\; x + y < 10$}
\end{enumerate}
¿Cuáles de estos enunciados son correctos? 
\begin{tasks}(3)
\task{Solo I}\task{Solo II}\task{Solo III}\task{Solo I y II}\task{Solo I y III}
\end{tasks}}
\end{itemize}
\end{document}